\subsection{应用：II. 有理整数的乘法}

引理2. 设 \(\mathbf{E}\) 是一个幺半群，并且 \(x\) 是 \(\mathbf{E}\) 的一个元素.

\begin{enumerate}
\def\labelenumi{(\roman{enumi})}
\item
  存在唯一的同态 \(f\) （从 \(\mathbf{N}\) 到 \(\mathbf{E}\) ）使得 \(f(1) = x\) 和 \(f(n) = \frac{n}{T} x\) （对所有 \(n \in \mathbf{N}\) ）.
\item
  若 \(x\) 是可逆的，存在唯一的同态 \(g\) （从 \(\mathbf{Z}\) 到 \(\mathbf{E}\) ）使得 \(g(1) = x\) 且 \(g\) 与 \(f\) 在 \(\mathbf{N}\) 上一致.
\end{enumerate}

记 \(f(n) = \top^n x\) （对所有 \(n \in \mathbf{N}\) ），则公式

\[
\overset {0} {\top} x = e \quad \text { and } \quad \left(\overset {m} {\top} x\right) \top \left(\overset {n} {\top} x\right) = \overset {m + n} {\top} x
\]

（第1号）表达了这样一个事实： \(f\) 是 \(\mathbf{N}\) 到 \(\mathbf{E}\) 的一个同态，且显然 \(f(1) = x\) 。若 \(f'\) 是 \(\mathbf{N}\) 到 \(\mathbf{E}\) 的一个同态，使得 \(f'(1) = x\) ，则 \(f = f'\) （根据第1节第4号命题1（iv））。

现在假设 \(x\) 是可逆的。根据第3号命题4的推论2， \(f(n) = \overset{n}{\top} x\) 对每个整数 \(n \geqslant 0\) 都是可逆的。根据构造， \(\mathbf{Z}\) 是 \(\mathbf{N}\) 的差群，因此（第4号定理1） \(f\) 唯一地延拓为一个同态 \(g\) （从 \(\mathbf{Z}\) 到 \(\mathbf{E}\) ）。若 \(g'\) 是 \(\mathbf{Z}\) 到 \(\mathbf{E}\) 的一个同态且 \(g'(1) = x\) ，则限制 \(f'\) （即 \(g'\) 到 \(\mathbf{N}\) 上的限制）是 \(\mathbf{N}\) 到 \(\mathbf{E}\) 的一个同态且 \(f'(1) = x\) 。因此 \(f' = f\) ，从而 \(g' = g\) .

我们将引理2应用于幺半群 E 为 Z 的情形；因此对每个整数 \(m \in Z\) 存在 Z 的一个自同态 \(f_{m}\) ，其特征是 \(f_{m}(1) = m\) 。若 m 属于 N，则映射 \(n \mapsto mn\) （从 N 到 N）是广群 N 的一个自同态（集合论，III，§3，第3号，命题5的推论）；因此 \(f_{m}(n) = mn\) （对所有 m, n \ensuremath{\in} N）。

\(\mathbf{N}\) 上的乘法因此可以扩张为 \(\mathbf{Z}\) 上的乘法，由公式 \(mn = f_m(n)\) （对于 \(m, n\) ，在 \(\mathbf{Z}\) 中）。我们将建立以下公式：

\[
x y = y x\tag{2}
\]

\[
(x y) z = x (y z)\tag{3}
\]

\[
x (y + z) = x y + x z\tag{4}
\]

\[
(x + y) z = x z + y z\tag{5}
\]

\[
0. x = x. 0 = 0\tag{6}
\]

\[
1. x = x. 1 = x\tag{7}
\]

\[
(- 1). x = x. (- 1) = - x\tag{8}
\]

对于 \(x, y, z\) （在 \(\mathbf{Z}\) 中）。（*换句话说， \(\mathbf{Z}\) 是一个交换环。*）公式 \(x(y + z) = xy + xz\) 和 \(x.0 = 0\) 表达了这一事实： \(f_x\) 是加法幺半群 \(\mathbf{Z}\) 的一个自同态，且 \(f_{x}(1) = x\) 可以写成 \(x.1 = x\) 。自同态 \(f_{x} \circ f_{y}\) （在 \(\mathbf{Z}\) 中）将 1 映到 \(xy\) ，故等于 \(f_{xy}\) ，由此得(3)。现在 \(f_{x}(-y) = -f_{x}(y)\) ，即 \(x(-y) = -xy\) ；类似地，自同态 \(y \mapsto -xy\) （在 \(\mathbf{Z}\) 中）将 1 映到 \(-x\) ，因此 \((-x).y = -xy\) ，从而

\[
(- x) (- y) = - (x (- y)) = - (- x y) = x y.
\]

对于 \(m, n\) （在 \(\mathbf{N}\) 中）, \(mn = nm\) （集合论，III，§3，第3号，命题5的推论），由此 \((-m).n = n(-m)\) 和 \((-m)(-n) = (-n)(-m)\) ；由于 \(\mathbf{Z} = \mathbf{N} \cup (-\mathbf{N})\) , \(xy = yx\) （对于 \(x, y\) ，在 \(\mathbf{Z}\) 中）；而由这个公式可知(5)由(4)推出，从而完成了公式(6)至(8)的证明。

